\documentclass[11pt]{article}
\usepackage[margin=1in]{geometry}
\usepackage[hidelinks]{hyperref}
\title{Upcoming Website Changes}
\author{Liam McGrath}
\date{Last updated 10 October 2026}
\begin{document}
\maketitle

Planned work, by the date it was requested. Each date lists the requested changes and how they are planned to be implemented. Finished items move to the Change Log.

\section{10 October 2026: open items}
\subsection{Requested changes}
\begin{itemize}
\item Confirm where Old Classics and My Hero Academia belong; they are under NanoFtheAbove for now.
\item Dynamic programming, trie and LeetCode Hard solutions, once they are in leetcode\_problems.
\end{itemize}
\subsection{Planned implementation}
\begin{itemize}
\item Moving a lyric folder is a one-line change to RAP\_PAGES in site\_content/catalog.py. A folder not yet placed appears under Future Projects, so nothing is dropped.
\item New solutions appear on their pages at the next build: dp or memo tables go to Dynamic Programming, and titles LeetCode rates Hard go to Hard Problems once they are added to HARD\_TITLES.
\end{itemize}

\section{Planned next: Series Hub subsections}
\subsection{Requested changes}
\begin{itemize}
\item Updates to each Series Hub destination: Videos, Analytics, History, Records and Creator Petition. Details will come with that prompt.
\item New basketball data for analysis in the style of a JimmyHighRoller video.
\end{itemize}
\subsection{Planned implementation}
\begin{itemize}
\item Extend the series data contract before adding new records, so earlier analysis stays reproducible.
\end{itemize}

\section{4 October 2026: open items}
\subsection{Requested changes}
\begin{itemize}
\item Owen's Creator Hub and Stock Monitoring links, once they are live.
\item Simulations in the style of Primer this month, as internal branches of the Simulations destination.
\item UFC and NBA predictions recorded against outcomes.
\item Swimming set times once a water-safe watch is bought: $5 \times 50$ m, $10 \times 50$ m and $5 \times 100$ m within set windows.
\item Science \& Engineering notes from the Gemini chat log; the contribution guide on the Mental Map explains the steps.
\item A database for change requests, once email requests become too many.
\end{itemize}
\subsection{Planned implementation}
\begin{itemize}
\item Owen's rows already exist and only need a URL each in creators.py.
\item The two predictions pages already exist; each needs a log of predictions made before events and the outcomes after.
\item The database plan follows below.
\end{itemize}

\section{Still open from 27 September 2026}
\subsection{Requested changes}
\begin{itemize}
\item Camille's creator section, once her site URL and written consent are in.
\item HTTPS on the custom domain: GitHub Pages still presents a certificate for the wrong name.
\end{itemize}
\subsection{Planned implementation}
\begin{itemize}
\item The HTTPS fix needs the repository owner in GitHub Pages settings: remove and re-add the custom domain, wait for the certificate, then turn on Enforce HTTPS.
\end{itemize}

\section{Database integration plan for change requests}
\label{database-plan}
Today, the Request a change form on the Mental Map drafts an email. When requests outgrow email, the same form can post to a small database instead. GitHub Pages only serves static files, so the database needs a separate service. The steps:

\begin{enumerate}
\item \textbf{Choose a service.} Either a hosted Postgres database with row-level security (for example Supabase), or a small serverless endpoint (for example a Cloudflare Worker) in front of a database (for example Cloudflare D1). The endpoint option keeps every secret off the website.
\item \textbf{Create the table.} Columns: id (generated), name (text, required, up to 120 characters), requested\_on (date, required), section, subsection and subsubsection (text, using the site's paths such as Science-and-Engineering/Physics/Optics), content (text, required, up to 20,000 characters), status (new, accepted, merged or declined, default new), reviewer\_note (text) and created\_at (timestamp).
\item \textbf{Validate on the server.} Accept only paths that exist in the site outline (assets/data/atlas-schema.js), require the subsubsection to sit inside the chosen subsection and section, enforce the lengths and strip control characters. Never trust the browser's checks alone.
\item \textbf{Lock down access.} The public may insert rows but never read, edit or delete them. With Supabase this is an insert-only policy for the anonymous role; with a Worker, the Worker holds the only key.
\item \textbf{Stop spam.} Add a bot check (for example Cloudflare Turnstile), a hidden honeypot field and a rate limit per address.
\item \textbf{Handle personal data.} Names are personal data: add a short privacy note to the form, keep requests only as long as needed and delete one on request.
\item \textbf{Switch the form.} In request-change.js, send the same fields as JSON to the endpoint, show a confirmation and keep the email draft as the fallback when the service is unavailable.
\item \textbf{Review requests.} Sort by status, export accepted requests, copy the content into mental\_map.tex under the right destination, then mark them merged and list them in the Change Log.
\item \textbf{Test.} Use a separate staging table, add tests for the validation rules, and try a malformed and an oversized request before going live.
\end{enumerate}

\end{document}
